\section{Normalized Walsh layers on tapes}\label{sec:tape-walsh}

The multiplication algorithm uses short Fourier transforms on several
address fields at once. Their factorization in Section~\ref{sec:integers}
requires a Walsh transform on one selected bit in each of several
fields. On one bit the transform is the normalized butterfly
\[
        (u,v)\longmapsto ((u+v)/\sqrt2,(u-v)/\sqrt2).
\]
Write \(H=2^{-1/2}\bigl(\begin{smallmatrix}1&1\\1&-1\end{smallmatrix}\bigr)\)
for its matrix. Given a set of bits in a \(k\)-bit address, a
\emph{selected-bit Walsh layer} applies \(H\) on those bits and the
identity on the rest. Fixing the other bits gives independent copies
of the same transform. The cost on tapes includes gathering those
copies, performing the transform at fixed precision, and returning
every record to the prescribed order.

\begin{lemma}[Selected-bit Walsh layers]\label{lem:tape-walsh}
Fix \(1/2<\rho<1\) and constants \(c,L,D>0\).
Let \(N=2^k\), \(k\ge1\), and specify
\(s\le ck/\log(2+k)\) distinct address bits. There is a complex
fixed-point format with fractional and signed integer widths equal to
fixed multiples of \(k\) in each real component. The multiples depend
only on \(\rho,c,L,D\). Given \(N\) records in lexicographic order
representing a vector \(x\) with \(\|x\|_2\le2^{Lk}\), a deterministic
fixed multitape machine computes the selected-bit Walsh layer to
absolute Euclidean error at most \(2^{-Dk}\). Its output has the
same format and order. It takes
\[
 O_{\rho,c,L,D}\!\left(Nk(1+k)^\rho(1+\log(2+k))\right)
\]
steps, including address changes and scalar preparation, and visits
\(O_{\rho,c,L,D}(Nk)\) cells. In the exact execution, the maps from the
input to any intermediate numerical array and from an additive error
in such an array to any later numerical array or output have operator
norm at most \(C(1+k)^B\) for fixed \(C,B\). The same bounds hold for
the adjoint layer.
\end{lemma}

\begin{proof}
Use Corollary~\ref{cor:bit-positions} to move the selected positions
to the low \(s\) bits. The other bits now index \(q=2^{k-s}\) spectator
rows, each holding a complete range of \(2^s\) addresses. The empty
selection is the identity. For odd \(s\) apply one butterfly directly
and transform the remaining \(s-1\) bits. Splitting into the two bit
classes aligns the operands for a scan of sums, differences and
scalings, followed by a merge. At the precision chosen below this
takes \(O(Nk\log(2+k))\) steps, including head returns. Bisection
with dyadic integer-square comparisons prepares \(1/\sqrt2\) to that
precision in \(k^{O(1)}\) steps. This settles \(s=1\). In every case
undo the initial bit-position permutation after the arithmetic.

\paragraph{A batch of transforms.}
Use the normalized binary construction in
\cite[Proposition~\ext{F}{prop:binary-template}, Equation~(\ext{F}{eq:kernel-template-saving})
and Section~\ext{F}{sec:prime-kernel}]{oaF}. It supplies fixed integers
\(A\ge3\), \(W\) and \(\sigma\). On \(W\) rows with \(Af\) address
bits per row, it implements a normalized transform of those bits by
\(\sigma\) child calls, each transforming one \(f\)-bit slot of one
work row on all its fibers. Here \(f\) is even. The strict saving is
\[
                            \sigma/W<A^\rho .
\]
We may duplicate the batch and all its calls to obtain \(W\ge2\).
Apart from the calls, the finite instruction list combines aligned
rows and changes their addresses. Pointwise scalar choices depend on
parities of sums of fixed functions of the \(f\) columns, each an
\(A\)-bit tuple with one bit from each slot. An address change applies
a fixed invertible binary matrix to every column. The scalars lie
in a fixed finite dyadic list because the normalized child transforms
already supply the normalization factors. The finite construction is
effective for rational \(\rho\). For any other admissible \(\rho\)
one can use a rational exponent strictly between \(1/2\) and \(\rho\).

Divide the even part of the selected width into pieces sized for this
recursion. Write \(\lfloor s/2\rfloor=\sum_{j=0}^J b_jA^j\), with
\(0\le b_j<A\). Allocate \(b_j\) consecutive pieces of width \(2A^j\)
and transform the pieces in succession. Append zero spectator rows to
make \(q'\), the padded leading row count, a multiple of \(W^J\).
Since \(W^J\le(1+s)^{O_\rho(1)}\) and \(q\ge2^{k/2}\) for large \(k\),
we have \(q'\le2q\). The bounded remaining widths admit direct
butterflies. Each full piece transform acts independently on all
spectator coordinates, so the added rows may be removed after the
last piece.

\paragraph{A recursive call on tapes.}
At a node with \(a=2A^j\) consecutive active bits, \(j>0\), the
leading row count \(r\) is divisible by \(W^j\). Set \(f=a/A\).
Split the leading rows by their residues modulo \(W\) onto \(W\)
streams. The entries at the same position on these streams supply the
aligned operands for one batch operation. Each stream has exactly
\(1/W\) of the parent's volume and leading count \(r/W\), divisible
by \(W^{j-1}\). Each child takes one work stream, designates one of
its \(f\)-bit slots as active, and leaves every other field as a full
spectator range. Figure~\ref{fig:walsh} shows this repeated layout.
The outputs merge by reversing the split. At \(a=2\) use the
four-point Walsh transform with coefficients \(\pm1/2\), splitting
and merging its four bit classes in a bounded number of scans.

\begin{figure}[!htb]
\centering
\begin{tikzpicture}[font=\small, x=1cm, y=1cm, line width=.52pt,
   descriptor/.style={anchor=east,align=right,font=\small},
   arrow/.style={draw=ink,-{Stealth[length=1.6mm]}}]
% Horizontal positions denote address fields, not their numeric widths.
\newcommand{\slotrow}[3]{%
  \fill[ink!8] (0,#1) rectangle (2.1,#1+.64);
  \fill[ink!8] (7.3,#1) rectangle (8.45,#1+.64);
  \draw[ink] (0,#1) rectangle (8.45,#1+.64);
  \foreach \cut in {1.05,2.1,3.4,4.7,6,7.3}
     \draw[ink] (\cut,#1)--(\cut,#1+.64);
  \node[font=\footnotesize] at (.525,#1+.32) {#2};
  \node[descriptor] at (-.23,#1+.32) {#3};}
\slotrow{2.35}{\(r\)}{Parent\\\(V\) bits}
\node[font=\footnotesize] at (1.575,2.67) {passive};
\node[font=\footnotesize] at (7.875,2.67) {passive};
\foreach \x/\txt in {2.75/1,4.05/2,5.35/3,6.65/4}
  \node at (\x,2.67) {\(\txt\)};
\draw[sea] (2.1,3.06)--(2.1,3.18)--(7.3,3.18)--(7.3,3.06);
\node[text=sea,above,font=\footnotesize] at (4.7,3.18)
  {\(A\) active slots of \(f\) bits each (\(A=4\) shown)};

\draw[arrow] (.525,2.26)--(.525,1.71)
  node[midway,right,font=\footnotesize,xshift=3pt]
  {split the \(r\) leading rows onto \(W\) streams};
\slotrow{1}{\(r/W\)}{One stream\\\(V/W\) bits}
\foreach \x in {2.45,3.75,5.05,6.35}
  \fill[rust!85] (\x,1.04) rectangle (\x+.12,1.60);
\node[font=\footnotesize,text=rust,anchor=west] at (6.77,1.92)
  {one column of four bits};
\draw[rust,line width=.45pt] (6.77,1.86)--(6.41,1.62);

\draw[arrow] (.525,.92)--(.525,.38)
  node[midway,right,font=\footnotesize,xshift=3pt]
  {recurse on one slot of this stream};
\slotrow{-.32}{\(r/W\)}{One call\\\(V/W\) bits}
\foreach \x in {2.1,4.7,6}
   \fill[ink!8] (\x+.025,-.30) rectangle (\x+1.275,.30);
\fill[sea!12] (3.425,-.30) rectangle (4.675,.30);
\draw[sea,line width=.8pt] (3.4,-.32) rectangle (4.7,.32);
\node[font=\footnotesize,text=sea] at (4.05,0) {\(f\) bits};
\end{tikzpicture}
\caption{One level of the Walsh recursion. The leading range counts
rows. The orange marks pick the same offset in each slot, so a
column change acts across slots. Each child transforms one whole slot
on every setting of its spectator fields. Other work streams remain
parked until the child returns.}
\label{fig:walsh}
\end{figure}


Let \(V\) be the bit volume of the node, including spectator ranges
and record widths. Binary elimination factors a column change into
a fixed number of slot exchanges and slot additions. An exchange uses
Theorem~\ref{thm:interchange}. For an addition, apply
Lemma~\ref{lem:packed-xor} with source and target slots of width \(f\)
and a third slot as auxiliary. All three ranges are complete, the
auxiliary contents are restored, and the leading range \(r/W\) may
have arbitrary length. These address changes together take
\[
 O_\rho\!\left(V(1+a^\rho)(1+\log(2+k))\right)+k^{O_\rho(1)}
\]
steps, including the narrow slots handled by direct bit operations
in Lemma~\ref{lem:packed-xor}. Pointwise passes take
\(O_\rho(V\log(2+k))+k^{O_\rho(1)}\) steps. To select a scalar, copy
the slots from the address counter onto separate tapes, scan all the
columns in parallel, accumulate the required parities, and rewind.
Even a scan of the full \(O(k)\)-bit address per record is paid for
by its width \(\Theta(k)\). There are only finitely many operands per
pointwise operation. Fixed dyadic scalings use short exact products
before rounding. A general product of two \(O(k)\)-bit numbers, as
in the direct butterfly, costs \(O(k\log(2+k))\) tape steps \cite{hvh}.

\paragraph{Work and visited space.}
Follow the depth-first schedule of \cite[Proposition~11]{Tape}.
Park the other streams and the parent descriptor on end-access stacks
while one child runs. On return, restore them to their marked work
areas. Copying, parking, erasing and rewinding add
\(O_\rho(V)+k^{O_\rho(1)}\) work per node. The child needs access only
to its current work areas and descriptors. Row counts, active widths,
and prefix, gap and suffix lengths fit in \(O_\rho(k)\)-bit descriptors
and all fixed-arity descriptor computations take \(k^{O_\rho(1)}\)
steps.

With the record width fixed throughout, divide time by volume and
write \(T(a)\) for the resulting bound at active width \(a\), omitting
the additive polynomial setup costs. The \(V/W\) volume of each child
gives
\[
 T(a)\le C_\rho(1+a^\rho)(1+\log(2+k))
                 +(\sigma/W)T(a/A).
\]
Relative to \(a^\rho\), successive levels have ratio
\(\sigma/(WA^\rho)<1\). Geometric summation with the four-point base
case yields \(T(a)=O_\rho(a^\rho(1+\log(2+k)))\). Also
\(\sum_j b_j(2A^j)^\rho=O_\rho(s^\rho)\). Each piece has constant
branching and depth \(O(\log(2+k))\), so all polynomial setup terms
together remain \(k^{O_\rho(1)}\), absorbed by \(N=2^k\). Including
both position permutations and the one-bit butterfly gives the
claimed time on a padded volume of \(O(Nk)\).

One node needs a fixed number of \(O(V)\)-bit streams for its scans
and for the sparse repair sort in Lemma~\ref{lem:packed-xor}.
Theorem~\ref{thm:interchange} also uses linear space. Along a single
descent, the parked volumes are bounded by a constant multiple of
\(V+V/W+V/W^2+\cdots\). Siblings and successive pieces reuse the
same tape regions. The additional descriptor and word workspace
along the stack is \(k^{O_\rho(1)}\), so the total number of visited
cells is \(O_{\rho,c,L,D}(Nk)\).

\paragraph{Precision.}
Give each concatenated set of live numerical arrays its Euclidean
norm. At one node there are a bounded number of operations on the
live arrays. Each has bounded norm when a completed child is counted
as one operation, since its transform on all fibers is unitary.
Copies and pointwise combinations have bounded norm from their
fixed coefficient and stream lists, and permutations are isometries.
For an input-to-intermediate estimate, descend only into the successive
calls containing the given array. An error originating inside a child
can propagate through its returned stream to the first common ancestor
of its operation and a later one, then down to the later array. At
most a constant norm factor is incurred per visited level. The depth
is logarithmic in \(k\), giving polynomial gain. If the two operations
occur in different pieces, all complete intervening pieces are
unitary on the padded array. Padding and restriction themselves have
norm at most one.

Round each real component after arithmetic to a grid of mesh \(2^{-p}\).
Choose \(1/\sqrt2\) to absolute error at most \(2^{-p}\) and store the
fixed dyadic coefficients exactly. Temporarily allow unbounded integer
parts. Rounding an array costs at most a constant times
\(\sqrt N\,2^{-p}\) in Euclidean norm. The approximate butterfly
coefficient introduces an additional error of at most a constant times
\(2^{-p}\) times the input norm. Thus while all preceding array norms
are at most \(2^{(L+2)k}\), a single operation incurs at most
\[
 C_\rho\,2^{-p}\bigl(\sqrt N+2^{(L+2)k}\bigr).
\]
Each array contains at most a fixed multiple of \(N\) entries. There
are \(k^{O_\rho(1)}\) array operations. Freeze their actual rounding
faults and propagate these vectors through the exact linear maps.
The polynomial gain bound then shows that \(p=\lceil C_{\rho,c,L,D}k\rceil\)
makes the total error at any array at most \(2^{-Dk}\), for a
sufficiently large constant. The corresponding ideal intermediate
norms are at most \(C(1+k)^B2^{Lk}\). For large \(k\) induction
therefore preserves the assumed norm bound, justifying signed integer
width \(O_{\rho,c,L,D}(k)\). Temporary exact sums and products before
rounding have the same order of width. Enlargement of the constant
width factors covers the bounded remaining lengths. Finally, the layer
is self-adjoint since \(H=H^*\).
\end{proof}
